\chapter{关键实现}\label{chap:impl}

本章结合项目实现，说明 QuickBoard 在“房间入口”“协作同步”“渲染优化”“公式识别与编辑”等方面的关键实现要点与工程取舍。

\section{房间号与入口归一化}
为了降低口述与输入成本，系统默认生成固定长度的纯数字房间号。首页加入房间支持直接输入房间号，也支持粘贴包含房间参数的链接或任意文本；前端在进入房间前对输入进行归一化，提取并校验房间号格式，从而减少用户在不同分享方式下的摩擦。

\subsection{最近房间与本地存储}
短会话协作具有“频繁进入同一房间或少量房间”的特点。为降低重复进入成本，前端将最近进入的房间号记录在浏览器本地存储，并在首页以列表形式展示，支持一键进入与移除记录。该设计不依赖账号体系即可实现轻量的“历史回访”，同时避免将房间记录写入服务端引入隐私与存储负担。

\subsection{房间生命周期与自动清理}
为了避免短会话结束后房间长期占用内存，后端实现空房间 TTL 机制：当房间人数为 0 且超过一定时间窗口后进行清理，并以固定间隔扫描回收。该机制与短会话场景的时间分布匹配，能够在不引入复杂持久化的前提下控制资源占用，并减少服务端长期运行时的状态堆积风险。

\section{服务端权威裁决与版本补包}
多人协作场景中，客户端操作可能乱序到达或重复投递。QuickBoard 将同步裁决放在服务端执行：客户端发送增量事件后，服务端根据对象级最后写入优先（LWW）策略决定是否接受并广播；同时维护房间版本号与增量日志。客户端加入房间时携带自身版本号，优先通过差量接口补齐缺失增量，必要时再回退到全量快照。该设计可以有效降低“刷新回滚”与“大包同步”的概率。

\subsection{字段级 LWW 与最小增量广播}
为避免“整对象覆盖”带来的误更新，服务端以字段粒度执行 LWW 裁决：仅当 incoming 的字段时间戳大于服务端权威状态中对应字段时间戳时才接受该字段更新。裁决后服务端广播被接受字段集合构成的最小增量，使房间内所有客户端接收一致的裁决结果，并减少网络与序列化开销。

\subsection{删除墓碑与对象复活防护}
在对象集合模型中，删除是不可逆语义。若删除仅以“本地移除对象”表达，则漏包客户端无法推断对象已被删除，旧修改事件晚到时还可能导致对象复活。QuickBoard 将删除表示为字段 \texttt{\_deleted=true} 并参与 LWW 合并：当 \texttt{\_deleted} 的时间戳更“新”时，服务端拒绝所有旧字段更新。快照对齐时携带 tombstones，使漏包客户端在下一次对齐时也能完成删除收敛。

\subsection{版本补包与退化为全量快照}
后端维护 \texttt{roomVersion} 与增量日志 \texttt{deltaLogs}：每当一次增量被服务端裁决并接受，\texttt{roomVersion} 递增并记录该版本对应的最小增量。客户端重连时携带 \texttt{clientVersion}：若日志覆盖缺失区间则下发 \texttt{sync-delta} 补齐；若缺失跨度超出日志覆盖能力则退化为 \texttt{sync-state} 权威快照。该退化策略使对齐语义清晰并优先保证正确性。

\section{渲染性能优化}
Fabric 在对象数量较多时容易出现重绘压力。项目在默认配置下开启视口裁剪与对象缓存，对静态对象进行缓存，选中对象时临时关闭缓存以保证编辑反馈。开发环境中提供基准入口用于对比不同渲染策略的帧耗时表现，从而在不引入复杂性能工具的前提下完成可复现的优化验证。

\section{Ghost Brush：过程预览降载}
过程预览点位属于高频瞬时数据。若逐点发送，不仅会造成消息风暴，也会挤占权威增量的带宽，间接放大一致性风险。QuickBoard 将过程预览与权威状态解耦：预览点位先入队，按固定周期批量发送；当队列过长时进行均匀抽样但保留尾部点位；单条消息设置最大点数并在必要时拆包。服务端仅转发过程预览，不写入权威状态与快照路径，从而避免过程噪声污染一致性链路。

\section{公式识别链路}
\subsection{截图与裁剪}
公式识别采用“白板选区截图”的方式采集输入。为避免坐标系不一致导致的空白截图，前端基于底层 Canvas 像素进行裁剪并生成图像数据，再发送给后端识别接口。该方式能够更稳定地反映选区内容，并减少不同缩放状态下的误差。

\subsection{OCR 预处理与结果清洗}
OCR 服务对输入图像进行尺寸对齐与内容裁剪，并在暗底场景下自动反色，必要时进行二值化生成多候选输入以提升识别成功率。对于识别结果中的噪声（如宏定义片段），服务端进行清洗并在清洗后为空时继续回退，最终对空结果统一返回错误码，便于前端给出一致的错误提示。

\subsection{后端代理、错误码与 traceId}
为降低前端与 OCR 服务的耦合，QuickBoard 采用后端代理调用：前端只请求后端 \texttt{/api/recognize-math}，由后端转发至 OCR 服务并统一超时控制与错误映射。后端对失败场景返回统一错误码（例如 \texttt{TIMEOUT}、\texttt{BAD\_IMAGE}、\texttt{EMPTY\_LATEX}、\texttt{RECOGNIZE\_FAILED}），并生成 \texttt{traceId} 随响应返回，使一次识别请求在前端、后端与 OCR 服务之间可被定位与追踪。该机制将“识别失败”从不可解释的黑盒结果转化为可定位的工程问题，便于论文中的失败案例分析与后续迭代。

\section{公式渲染、确认插入与二次编辑}
识别得到的 LaTeX 并不直接以纯文本显示在画布上，而是通过 KaTeX 渲染为可读的数学公式。交互上，系统采用“识别后先确认再插入”的流程，避免误识别内容污染画布；插入后支持双击打开编辑弹层并实时预览渲染结果，确认后更新公式对象。为保证选框与实际渲染区域一致，公式对象在画布中的命中框会根据渲染尺寸进行补边与校正，从而提升选择与拖拽的可用性。

\subsection{超时控制与交互不中断}
公式识别涉及网络请求与模型推理，耗时具有不确定性。为避免识别请求长时间阻塞交互，前端在请求侧加入超时控制并支持中止，使用户在网络波动或服务端卡顿时仍可继续使用白板其他功能。后端同样设置请求超时并将失败映射为统一错误码，使“超时”在系统内部具有一致语义，便于前端提示与论文评估。

\subsection{对象序列化与一致显示}
为保证公式对象在多端显示一致，系统需要在对象序列化中携带与公式渲染相关的必要字段（例如 LaTeX 字符串与渲染尺寸信息）。在同步与对齐过程中，远端客户端根据这些字段重建对象，并使用同一渲染策略生成显示结果。相较于传输位图，传输 LaTeX 文本具有可编辑、体积小与跨分辨率稳定的优势；同时，命中框与边距的校正使选择、拖拽与缩放的交互体验更接近普通图形对象。
